\chapter{Problem Definition}
\noindent
Despite significant advances in both automatic speech recognition (ASR) and natural language summarization, several critical gaps remain in the meeting summarization domain.

\section{Modality Isolation}
Most existing systems process text, audio, and visual streams independently, failing to capture cross-modal correlations that indicate segment importance. For instance, a speaker switching to a specific slide while increasing their vocal pitch usually signals high importance, a cue missed by unimodal systems.

\section{Temporal Discontinuity}
Current approaches treat each meeting as an isolated event. This leads to the loss of valuable context regarding ongoing discussions, recurring topics, and evolving decisions that span across a meeting series.

\section{Role Agnosticism}
Generic summaries do not address the diverse information needs of different stakeholders. A product manager requires feature-centric views, while an engineer requires technical specifications. 

\section{Static Fusion Weights}
Rule-based fusion approaches use fixed weights for combining modalities, which are unable to adapt to varying meeting contexts where different signals (visual vs. audio) may carry varying importance.

The core problems addressed in this work are:
\begin{itemize}
    \item \textbf{Loss of Inter-Meeting Context:} Standard meeting summarizers treat each meeting in isolation. Topics, outstanding questions, and evolving decisions lose their link across different meetings, leading to disjointed knowledge.
    \item \textbf{Ignored Multimodal Context:} Relying solely on audio transcriptions misses critical visual information—such as data displayed on shared screens or presentations—which is essential for accurately capturing the meeting's complete technical context.
    \item \textbf{The "Action" Gap:} Summarization alone does not execute tasks. Even a perfect summary requires manual effort to translate decisions into tangible workflow updates, tickets, or calendar events.
\end{itemize}
RoME aims to solve these issues by marrying multimodal extraction with a persistent temporal knowledge graph and an autonomous, tool-equipped AI agent.
